\documentclass[11pt]{article}

% ---------------------------------------------------------------------------
% Portable preamble: compiles with pdfLaTeX or XeLaTeX, no project macros.
% ---------------------------------------------------------------------------
\usepackage{amsmath,amsthm,amssymb}
\usepackage[margin=1in]{geometry}
\usepackage[colorlinks=true,linkcolor=blue,urlcolor=blue,citecolor=blue]{hyperref}
\usepackage{mathtools}
\usepackage{enumitem}
\usepackage{booktabs}
\usepackage{longtable}
\usepackage[round,authoryear]{natbib}

\theoremstyle{plain}
\newtheorem{theorem}{Theorem}[section]
\newtheorem{lemma}[theorem]{Lemma}
\newtheorem{proposition}[theorem]{Proposition}
\newtheorem{corollary}[theorem]{Corollary}
\theoremstyle{definition}
\newtheorem{assumption}{Assumption}
\newtheorem{definition}[theorem]{Definition}
\theoremstyle{remark}
\newtheorem{remark}[theorem]{Remark}
\newtheorem{example}[theorem]{Example}

\renewcommand{\theassumption}{A\arabic{assumption}}

\newcommand{\E}{\mathbb{E}}
\newcommand{\PP}{\mathbb{P}}
\newcommand{\Pn}{\mathbb{P}_n}
\newcommand{\R}{\mathbb{R}}
\newcommand{\cX}{\mathcal{X}}
\newcommand{\cT}{\mathcal{T}}
\newcommand{\cS}{\mathcal{S}}
\newcommand{\cP}{\mathcal{P}}
\newcommand{\cL}{\mathcal{L}}
\newcommand{\ez}{e_0}
\newcommand{\mz}{\mu_0}
\newcommand{\wz}{w_0}
\newcommand{\ehat}{\widehat{e}}
\newcommand{\mhat}{\widehat{\mu}}
\newcommand{\what}{\widehat{w}}
\newcommand{\that}{\widehat{\tau}}
\newcommand{\thetahat}{\widehat{\theta}}
\newcommand{\thetatil}{\widetilde{\theta}}
\newcommand{\pihat}{\widehat{\pi}}
\newcommand{\Vhat}{\widehat{V}}
\newcommand{\Dw}{D_w}
\newcommand{\Dm}{D_\mu}
\newcommand{\Ltwo}[1]{\lVert #1 \rVert_{2}}
\newcommand{\Lw}[1]{\lVert #1 \rVert_{2,w}}
\newcommand{\Linf}[1]{\lVert #1 \rVert_{\infty}}
\DeclareMathOperator*{\argmin}{arg\,min}
\DeclareMathOperator*{\argmax}{arg\,max}
\DeclareMathOperator{\Var}{Var}
\DeclareMathOperator{\plim}{plim}
\DeclarePairedDelimiter{\abs}{\lvert}{\rvert}
\DeclarePairedDelimiter{\card}{\lvert}{\rvert}

\title{Interpretable Causal Inference with Optimal Decision Trees\\[2pt]
\large Theory}
\author{}
\date{}

\begin{document}
\maketitle

\begin{abstract}
We develop inference for the average treatment effect on the treated when both
nuisance functions --- the propensity score and the control-outcome regression
--- are estimated by globally optimal decision trees. The covariate space is
finite with cardinality fixed in the sample size, and the leaf budget is fixed,
so the class of candidate tree partitions is a fixed finite set. Our main result
is that a single tree per nuisance, with both the partition and the leaf values
fitted on all $n$ observations, yields $\sqrt n$-consistent, asymptotically
normal, and locally asymptotically efficient estimation of the target, with
consistent variance estimation by the ordinary empirical influence-function
plug-in. No sample splitting, cross-fitting, or tail condition is used. Two
structural facts drive the argument: the induced estimand is invariant across
every partition rich enough to represent the truth, so data-adaptive selection
of the tree structure cannot move the target; and a positive population risk
margin is a derived consequence of exact global optimisation over a finite
class, rather than an assumption. We further give a cross-fitted estimator that
dispenses with the structural sparsity condition, and confidence intervals that
remain valid, at a stated and estimable price in width, when sparsity fails.
\end{abstract}

\tableofcontents

% ===========================================================================
\section{Scope and organisation}\label{sec:scope}

This document is the canonical statement of the theory, together with its
proofs. Section~\ref{sec:setup} fixes notation and collects every standing
assumption; assumptions are defined there and nowhere else, and are invoked by
label throughout. Section~\ref{sec:eif} records the efficient influence function
and the exact second-order remainder. Section~\ref{sec:structure} develops the
structure of the tree class: sufficiency, invariance, the margin, and selection
consistency. Section~\ref{sec:main} contains the principal result, full-sample
inference with two optimal trees. Section~\ref{sec:crossfit} treats the
cross-fitted estimator, which requires no sparsity condition and serves as the
anchor for Section~\ref{sec:honest}, where inference remains valid but
inefficient when sparsity fails. Section~\ref{sec:discussion} discusses the
results and their relation to adjacent literatures.

Proofs are collected in Appendices~\ref{app:eif}--\ref{app:honest}, keyed to
result numbers; auxiliary lemmas are in Appendix~\ref{app:aux} and the
counterexamples are worked in full in Appendix~\ref{app:counterex}. A proof is
given inline only where the proof is the substance of the result and occupies a
few lines. Every numbered statement in this document has a proof. Claims that
are not proved appear as unnumbered prose in Section~\ref{sec:discussion}.

Two bodies of theory are deliberately excluded and are developed in a companion
paper: approximation and oracle theory for optimal trees on continuous
covariates, where the leaf budget must grow with $n$ and entropy conditions
become active; and Rashomon sets, together with the structure-intersection
results that support them. Neither is used by any result below.

Appendix~\ref{app:crosswalk} maps each numbered result to its destination in the
submitted paper.

% ===========================================================================
\section{Setup and standing assumptions}\label{sec:setup}

\subsection{Data, target, and score}

Let $O=(X,A,Y)$ with covariate $X$, binary treatment $A\in\{0,1\}$, and outcome
$Y\in\R$, distributed according to $P$. Let $Y(1),Y(0)$ denote the potential
outcomes. Write
\[
  \pi=\PP(A=1),\qquad
  \ez(x)=\PP(A=1\mid X=x),\qquad
  \mz(x)=\E[Y\mid A=0,X=x],
\]
and let $m_1(x)=\E[Y\mid A=1,X=x]$ and $\sigma_a^2(x)=\Var(Y\mid A=a,X=x)$. The
target is the average treatment effect on the treated,
\[
  \theta_0=\E[Y(1)-Y(0)\mid A=1].
\]
The ATT weight is $\wz=\ez/(1-\ez)$, and for a candidate propensity $e$ we write
$w=e/(1-e)$. A candidate nuisance pair is $\eta=(e,\mu)$, with $\eta_0=(\ez,\mz)$
the truth. It is convenient to index the two nuisances by $j\in\{e,\mu\}$ and to
write $\gamma_{0,e}=\ez$, $\gamma_{0,\mu}=\mz$.

We write $\Pn$ for the empirical measure, $\Pn f=n^{-1}\sum_{i\le n}f(O_i)$, and
$\pihat=\Pn A$. Two norms on functions of $x$ are used throughout:
\[
  \Ltwo f^2=\E\bigl[f(X)^2\bigr],
  \qquad
  \Lw f^2=\E\bigl[\{1-\ez(X)\}f(X)^2\bigr],
\]
the second being the control-weighted norm. Convergence in distribution is
denoted $\rightsquigarrow$, and $z_{1-\alpha/2}$ is the standard normal
quantile.

\subsection{Covariate space and tree partitions}

The covariate space is a finite product,
\[
  \cX=\prod_{j=1}^d\cX_j,\qquad \card{\cX_j}=J_j,\qquad
  M=\card{\cX}=\prod_{j\le d}J_j ,
\]
whose elements we call \emph{cells}, with cell probabilities $p_x=\PP(X=x)$ and
$p_{\min}=\min_{x}p_x$.

A \emph{tree partition} $\tau$ of $\cX$ is a partition induced by recursive
axis-aligned binary splits on the coordinates of $\cX$; its blocks are
\emph{leaves}, $\card\tau$ is its leaf count, and $\ell_\tau(x)$ is the leaf
containing $x$. For a leaf budget $\bar L$,
\[
  \cT_{\bar L}=\bigl\{\tau:\tau\text{ a tree partition of }\cX,\
  \card\tau\le\bar L\bigr\}.
\]
For a strictly positive weight $\nu$ on $\cX$, let $\Pi^\nu_\tau$ denote the
$\nu$-weighted projection onto functions constant on the leaves of $\tau$; we
write $\Pi_\tau$ for $\Pi^{p}_\tau$ and use the control-weighted projection,
$\nu_x=p_x\{1-\ez(x)\}$, for the outcome nuisance. Given $\tau$, the
\emph{full-sample leaf refit} sets, on each leaf $\ell$,
\begin{equation}\label{eq:refit}
  \ehat_\tau\restriction_\ell
  =\frac{\sum_{i:X_i\in\ell}A_i}{\card{\{i:X_i\in\ell\}}},
  \qquad
  \mhat_\tau\restriction_\ell
  =\frac{\sum_{i:X_i\in\ell,\,A_i=0}Y_i}
        {\card{\{i:X_i\in\ell,\,A_i=0\}}} .
\end{equation}

Structure is selected by penalised empirical risk minimisation: for a penalty
$\lambda_n>0$,
\begin{equation}\label{eq:select}
  \that_j\in\argmin_{\tau\in\cT_{\bar L}}
  \bigl\{R^{(j)}_n(\tau)+\lambda_n\card\tau\bigr\},
  \qquad j\in\{e,\mu\},
\end{equation}
where $R^{(j)}_n$ is the empirical risk of the leaf-refit fit for nuisance $j$
and $R^{(j)}$ its population counterpart. We write $\ehat=\ehat_{\that_e}$,
$\mhat=\mhat_{\that_\mu}$, $\what=\ehat/(1-\ehat)$, and
$L_e=\card{\that_e}$, $L_\mu=\card{\that_\mu}$.

\subsection{Errors of the fitted nuisances}

Two distinct error notions appear, and they are not interchangeable. The
\emph{realised} errors are properties of the nuisances actually reported,
\begin{equation}\label{eq:realised}
  \Dw=\Lw{\what-\wz},\qquad \Dm=\Lw{\mhat-\mz},
\end{equation}
whereas the \emph{approximation} errors are the best available within the
budget,
\begin{equation}\label{eq:approx}
  \delta_e=\min_{\tau\in\cT_{\bar L}}\Ltwo{\Pi_\tau\ez-\ez},
  \qquad
  \delta_\mu=\min_{\tau\in\cT_{\bar L}}\Lw{\Pi^{\nu}_\tau\mz-\mz}.
\end{equation}
The realised error is a random quantity attached to a particular fit; the
approximation error is a deterministic functional of $P$ and $\cT_{\bar L}$.
Results stated on one do not transfer to the other.

\subsection{Standing assumptions}

Assumptions~\ref{ass:iid}--\ref{ass:finite} restrict the sampling model,
\ref{ass:construct}--\ref{ass:global} the estimator, and
\ref{ass:sparsity}--\ref{ass:pseudo} the structure of the problem. Each result
below states which it uses. Assumptions \ref{ass:margin-floor} and
\ref{ass:pseudo} are local to Sections~\ref{sec:uniformity}
and~\ref{sec:realised} respectively and are never in force elsewhere.

\begin{assumption}[Sampling]\label{ass:iid}
$O_1,\dots,O_n$ are independent and identically distributed copies of
$O=(X,A,Y)\sim P$.
\end{assumption}

\begin{assumption}[Identification and one-sided positivity]\label{ass:causal}
Consistency, $Y=AY(1)+(1-A)Y(0)$; ignorability,
$\{Y(0),Y(1)\}\perp A\mid X$; $\pi>0$; and there exists $c\in(0,1)$ with
\[
  1-\ez(x)\ \ge\ c \qquad\text{for every }x\in\cX .
\]
\end{assumption}

\noindent Positivity is required on one side only. Identification of $\theta_0$
uses no uniform constant: the treated mean $\E[Y(1)\mid A=1]$ is available from
treated units alone, and cells with $\ez(x)=0$ receive weight zero in
$\theta_0=\pi^{-1}\E[\ez\,(m_1-\mz)]$. Inference uses the uniform bound
$1-\ez\ge c$ and nothing further: it is equivalent to boundedness of $\wz$, it
is what makes the second-order remainder of Section~\ref{sec:eif} finite, and it
supplies the positive control mass $p_x\{1-\ez(x)\}\ge cp_x$ that the outcome
margin requires. No lower bound on $\ez$ is used anywhere, and $\ez(x)=0$ is
harmless.

\begin{assumption}[Moments and non-degeneracy]\label{ass:moments}
$\E[Y^2\mid A=a,X=x]\le C_Y<\infty$ for all $a\in\{0,1\}$ and $x\in\cX$, and
$V>0$, where $V$ is given in~\eqref{eq:vdecomp}.
\end{assumption}

\begin{assumption}[Finite covariate space, fixed]\label{ass:finite}
$\cX$ is finite with $d$, each $J_j$, and hence $M=\card\cX$ fixed in $n$; and
$p_x>0$ for every $x\in\cX$.
\end{assumption}

\begin{assumption}[Estimator construction]\label{ass:construct}
The nuisances are fitted as follows.
\begin{enumerate}[label=(\alph*),nosep]
  \item \emph{Loss.} Both nuisances are fitted by squared error, so that excess
  risk equals squared $L_2$ distance to the truth.
  \item \emph{Leaf refit.} Leaf values are given by~\eqref{eq:refit}: the
  within-leaf treated fraction for $\ehat$, the within-leaf control mean for
  $\mhat$.
  \item \emph{Clipping.} Fitted propensities are clipped so that
  $1-\ehat(x)\ge c$ for every $x$, with $c$ the constant of
  Assumption~\ref{ass:causal}.
  \item \emph{Leaf mass.} Every leaf of $\that_e$ contains at least $m_n$
  observations and every leaf of $\that_\mu$ at least $m_n$ control
  observations, with $\bar Lm_n\lesssim n$. On the null event that a leaf of
  $\that_\mu$ contains no controls, $\mhat$ takes the parent leaf's control
  mean.
\end{enumerate}
\end{assumption}

\noindent Three comments on Assumption~\ref{ass:construct}. Under (a) the
excess risk of a candidate $f$ for the propensity is
$\E[(A-f)^2\mid X]-\E[(A-\ez)^2\mid X]=(\ez-f)^2$ exactly, for any
$\ez\in[0,1]$, with no positivity; the loss-to-norm constant is $1$. The
within-leaf treated fraction minimises both squared error and log-loss over
constants, so (b) is invariant to that choice of loss and the loss affects only
the structure search~\eqref{eq:select}. Clipping in (c) is applied on one side
because Assumption~\ref{ass:causal} places $\ez$ in $[0,1-c]$, on which
one-sided clipping is the Euclidean projection and therefore non-expansive:
$\abs{\mathrm{clip}(\ehat)-\ez}\le\abs{\ehat-\ez}$ pointwise, so every $L_2$
bound is preserved. Two-sided clipping forfeits this wherever $\ez<c$. The clip
level must equal $c$: clipping at $c'>c$ changes the influence function on cells
with $\ez>1-c'$, so that $\thetahat$ remains consistent by double robustness
while $V$ and every coverage statement below fail.

\begin{assumption}[Fixed leaf budget]\label{ass:budget}
$\bar L$ is a constant, fixed in $n$, with $\bar L\le M$; consequently
$\cT_{\bar L}$ is a fixed finite set and $\card{\cT_{\bar L}}$ does not depend
on $n$.
\end{assumption}

\noindent Assumptions~\ref{ass:finite} and~\ref{ass:budget} together are what
make the analysis below elementary:
$\log\card{\cT_{\bar L}}\asymp\bar L\log(dJ)=O(1)$, so a union bound over the
candidate class costs a constant, no exponential tail inequality is needed at
any point, and Assumption~\ref{ass:moments} suffices with no sub-Gaussian
condition. The results of this document should not be motivated as escaping an
entropy cost, because under \ref{ass:finite}--\ref{ass:budget} no such cost
exists.

\begin{assumption}[Exact global optimisation]\label{ass:global}
For each $j$, $\that_j$ attains the global minimum in~\eqref{eq:select} over
$\cT_{\bar L}$, and $\lambda_n\to0$.
\end{assumption}

\begin{definition}[Sufficient class]\label{def:sufficient}
For $j\in\{e,\mu\}$,
\[
  \cS_j=\bigl\{\tau\in\cT_{\bar L}:
  \Pi^\nu_\tau\gamma_{0,j}=\gamma_{0,j}\ \ P\text{-a.s.}\bigr\},
\]
equivalently the set of $\tau\in\cT_{\bar L}$ on whose leaves $\gamma_{0,j}$ is
$P$-a.s.\ constant. The characterisation is independent of the strictly positive
weight $\nu$.
\end{definition}

\begin{assumption}[Structural sparsity]\label{ass:sparsity}
$\cS_e\ne\emptyset$ and $\cS_\mu\ne\emptyset$; equivalently
$\delta_e=\delta_\mu=0$.
\end{assumption}

\noindent Assumption~\ref{ass:sparsity} is a sparsity condition, not a
correct-specification condition. On a finite covariate space every function is
piecewise constant, and the partition into singletons is tree-realisable using
$M-1$ splits, so at $\bar L=M$ the assumption holds automatically and carries no
content. Its content at $\bar L<M$ is that a \emph{small} tree suffices, that
is, that the nuisances are structurally sparse and hence displayable.

The cost of Assumption~\ref{ass:sparsity} is not symmetric in the type of
structure present. With $s$ active coordinates in a region, a tree needs at most
$\prod_{j\le s}J_j$ leaves, and far fewer when dependence is
\emph{hierarchical}, in the sense that the active coordinate set differs across
regions: as few as $s+1$ leaves. The assumption is expensive precisely for
\emph{additive} structure, where a leaf is required for each configuration of
the active coordinates, and cheap precisely for region-varying active sets. That
asymmetry, rather than flexibility, is the reason to use trees here; sparsity is
local, so $\gamma_{0,j}$ may depend on all $d$ coordinates globally while
remaining cheap to represent.

\begin{assumption}[Margin floor]\label{ass:margin-floor}
$\min_{j\in\{e,\mu\}}\Delta_j\ge\Delta_{\min}>0$, with $\Delta_j$ as in
Lemma~\ref{lem:margin}.
\end{assumption}

\begin{assumption}[Pseudo-true margin]\label{ass:pseudo}
There is a unique pair $(\tau_e^\dagger,\tau_\mu^\dagger)\in\cT_{\bar L}^2$
minimising the population risks $R^{(e)}$ and $R^{(\mu)}$ over
$\cT_{\bar L}$, and
\[
  \Delta^\dagger_j:=\min_{\tau\in\cT_{\bar L}\setminus\{\tau_j^\dagger\}}
  \bigl\{R^{(j)}(\tau)-R^{(j)}(\tau^\dagger_j)\bigr\}>0,
  \qquad j\in\{e,\mu\}.
\]
Write $\eta_*=(e_*,\mu_*)$ for the corresponding population leaf-refit
nuisances, $w_*=e_*/(1-e_*)$, and $V_*$ for the asymptotic variance of the
estimator at $\eta_*$.
\end{assumption}

\noindent Assumption~\ref{ass:pseudo} is strictly weaker than
Assumption~\ref{ass:sparsity}: it requires a unique population-risk minimiser
within the budget and a gap to the nearest competitor, but not that the
minimiser represent the truth. Under \ref{ass:sparsity} it holds with
$\tau^\dagger_j$ the finest member of $\cS_j$ and $\eta_*=\eta_0$.

% ===========================================================================
\section{The influence function and the exact remainder}\label{sec:eif}

\subsection{Score and estimator}

The efficient influence function for $\theta_0$ in the nonparametric model is
\citep{hahn1998}
\begin{equation}\label{eq:score}
  \psi(O;\theta,\eta)
  =\frac1\pi\Bigl[A\bigl\{Y-\mu(X)-\theta\bigr\}
   -(1-A)\,w(X)\bigl\{Y-\mu(X)\bigr\}\Bigr],
  \qquad w=\frac e{1-e},
\end{equation}
and we write $\psi_0=\psi(\cdot;\theta_0,\eta_0)$. Solving
$\Pn\psi(\cdot;\theta,\widehat\eta)=0$ gives the estimator
\begin{equation}\label{eq:estimator}
  \thetahat=\frac1{n\pihat}\sum_{i=1}^n
  \Bigl[A_i\bigl\{Y_i-\mhat(X_i)\bigr\}
   -(1-A_i)\,\what(X_i)\bigl\{Y_i-\mhat(X_i)\bigr\}\Bigr].
\end{equation}
The control-arm term carries a minus sign, as solving $\Pn\psi=0$ requires. The
treated share is not substituted into~\eqref{eq:estimator} but emerges from the
coefficient of $\theta$ in~\eqref{eq:score}: with $\phi:=\pi\psi$,
\begin{equation}\label{eq:exact-pi}
  \pihat\,(\thetahat-\theta_0)
  =\Pn\phi(\cdot;\theta_0,\widehat\eta)
\end{equation}
holds exactly, with no expansion, so that no further correction for the
estimation of $\pi$ is required.

\subsection{The second-order remainder}

\begin{proposition}[Exact bilinear remainder]\label{prop:bilinear}
Under Assumptions~\ref{ass:iid}--\ref{ass:causal}, for any $\eta=(e,\mu)$ with
$1-e\ge c$,
\begin{equation}\label{eq:bilinear}
  \E\bigl[\psi(O;\theta_0,\eta)\bigr]
  =\frac1\pi\,\E\Bigl[\bigl\{1-\ez(X)\bigr\}
   \bigl\{\wz(X)-w(X)\bigr\}\bigl\{\mz(X)-\mu(X)\bigr\}\Bigr].
\end{equation}
\end{proposition}

Equation~\eqref{eq:bilinear} is the classical doubly robust computation
\citep{robins1995,vanderlaan2003,kennedy2024}; it is recorded here because the
exact form of its weighting determines the constants in
Lemma~\ref{lem:biasbound} and Corollary~\ref{cor:convert}, which differ, and
because Section~\ref{sec:honest} requires it at the fitted rather than the
population nuisances.

\begin{lemma}[Bias bound in the control-weighted norm]\label{lem:biasbound}
Under Assumptions~\ref{ass:iid}--\ref{ass:causal} and
Assumption~\ref{ass:construct}(c),
\[
  \bigl\lvert\E[\psi(O;\theta_0,\widehat\eta)]\bigr\rvert
  \ \le\ \frac1\pi\,\Dw\,\Dm .
\]
\end{lemma}

\noindent The constant is exactly $1/\pi$: Cauchy--Schwarz with respect to
$\{1-\ez\}\,dP$ places $\sqrt{1-\ez}$ on each factor, and no power of $c$
appears. Converting to the conventional unweighted propensity error is not free.

\begin{corollary}[Unweighted form]\label{cor:convert}
Under the conditions of Lemma~\ref{lem:biasbound},
\begin{equation}\label{eq:convert}
  \bigl\lvert\E[\psi(O;\theta_0,\widehat\eta)]\bigr\rvert
  \ \le\ \frac1{\pi c^{3/2}}\,\Ltwo{\ehat-\ez}\,\Lw{\mhat-\mz}.
\end{equation}
\end{corollary}

\noindent The asymmetry in~\eqref{eq:convert} --- unweighted for the propensity,
control-weighted for the outcome --- is not a choice. One full factor of
$1-\ez$ is consumed by the identity
$\{1-\ez\}\abs{\wz-w}=\abs{\ez-e}/(1-e)$, which leaves the outcome factor to be
converted at cost $c^{-1/2}$. Since $c\le1-\ez\le1$ gives
$c\Ltwo f^2\le\Lw f^2\le\Ltwo f^2$, the norm-equivalence constant is $1/c$ for
squared norms and $c^{-1/2}$ for norms. Form~\eqref{eq:convert} is what the rate
condition of Section~\ref{sec:crossfit} uses; Lemma~\ref{lem:biasbound} is what
Section~\ref{sec:honest} uses.

\subsection{The efficiency bound}

\begin{lemma}[Tangent space and the variance bound]\label{lem:tangent}
Let $\cP$ be the set of laws of $(X,A,Y)$ satisfying
Assumptions~\ref{ass:causal}--\ref{ass:finite}. Then the tangent space of $\cP$
at any $P_0\in\cP$ is infinite-dimensional; $\theta$ is pathwise differentiable
at $P_0$ relative to $\cP$ with canonical gradient $\psi_0$; and the
semiparametric efficiency bound for $\theta_0$ is $V=\E[\psi_0^2]$, given
explicitly by
\begin{equation}\label{eq:vdecomp}
  V=\pi^{-2}\,\E\Bigl[\ez\sigma_1^2
   +\frac{\ez^2\sigma_0^2}{1-\ez}
   +\ez\,(m_1-\mz-\theta_0)^2\Bigr].
\end{equation}
\end{lemma}

\noindent Finiteness of $\cX$ does not make the problem parametric. With $M$
fixed, $P$ is determined by the cell probabilities $\{p_x\}$ and the propensity
$\{\ez(x)\}$ \emph{together with} the conditional laws of $Y$ given $(A,X)$,
which Assumption~\ref{ass:moments} leaves unrestricted. The tangent space
accordingly contains all square-integrable $b(y,a,x)$ with
$\E[b\mid A,X]=0$, and $V$ in~\eqref{eq:vdecomp} is a genuine semiparametric
bound rather than a parametric information bound.

% ===========================================================================
\section{Structure of the tree class}\label{sec:structure}

Throughout this section $j\in\{e,\mu\}$ and $\cS_j$ is as in
Definition~\ref{def:sufficient}. The order of results is deliberate:
invariance of the estimand (Lemma~\ref{lem:invariance}) precedes the margin
(Lemma~\ref{lem:margin}), which precedes the failure of the margin under greedy
splitting (Proposition~\ref{prop:greedy}), which precedes selection consistency
(Lemma~\ref{lem:selection}).

\subsection{Sufficiency is intrinsic, not a refinement condition}

Definition~\ref{def:sufficient} characterises sufficiency intrinsically, by the
requirement that the projection of the truth onto $\tau$ return the truth. It is
tempting to characterise it instead by refinement of a minimal partition on
which $\gamma_{0,j}$ is constant, and to read the sparsity condition as a
comparison of $\bar L$ with the cardinality of that partition. The two are not
equivalent, because a minimal partition into axis-aligned \emph{boxes} need not
be realisable as a \emph{tree}.

\begin{proposition}[Minimal box partitions need not be tree-realisable]
\label{prop:pinwheel}
There exist a finite product space $\cX$ with all cells of positive
probability, a function $\gamma_0$ on $\cX$, and a budget $\bar L$ equal to the
cardinality of a minimal partition of $\cX$ into axis-aligned boxes on which
$\gamma_0$ is constant, such that no tree partition in $\cT_{\bar L}$ has
$\gamma_0$ constant on its leaves. That is, $\cS=\emptyset$ while
$\bar L$ equals the minimal box complexity of $\gamma_0$.
\end{proposition}

\noindent A five-box witness on $\cX=\{0,1,2\}^2$ requiring six leaves is
constructed in Appendix~\ref{app:counterex}. The consequence is that the
refinement formulation can render $\cS_j$ empty exactly when the budget appears
to be adequate, in which case Lemma~\ref{lem:selection} has probability zero and
Theorem~\ref{thm:main} has no valid route. Definition~\ref{def:sufficient}
avoids the distinction between boxes and trees, and also avoids a support
qualifier, since it is stated $P$-a.s.

\subsection{Invariance of the estimand}

\begin{lemma}[Invariance over the sufficient class]\label{lem:invariance}
Under Assumptions~\ref{ass:iid}--\ref{ass:causal}, \ref{ass:finite}, and
\ref{ass:sparsity}, for every $(\tau_e,\tau_\mu)\in\cS_e\times\cS_\mu$ the
population leaf-refit nuisances $\eta_\tau=(e_{\tau_e},\mu_{\tau_\mu})$ satisfy
$\theta(\eta_\tau)=\theta_0$, where $\theta(\eta)$ solves
$\E[\psi(O;\theta,\eta)]=0$.
\end{lemma}

Lemma~\ref{lem:invariance} is the load-bearing structural fact. The target of
inference does not move as the selected structure ranges over the sufficient
class; consequently $\that_j$ need never recover a particular partition, and no
\emph{rate} for structure selection is required. Membership $\that_j\in\cS_j$
with probability tending to one suffices. It also disposes of a distinction that
would otherwise require attention: $\mhat$ estimates the control-weighted
within-leaf mean of $\mz$ rather than the unweighted mean, but by
Definition~\ref{def:sufficient} the two coincide on the sufficient class.

The relation between Lemma~\ref{lem:invariance} and the impossibility results
for post-selection inference \citep{leeb2005} is direct. Those results concern
inference on a parameter whose meaning changes with the selected model. Here it
does not: the estimand is constant on $\cS_e\times\cS_\mu$.

\subsection{The margin is derived}

\begin{lemma}[Positive population margin]\label{lem:margin}
Under Assumptions~\ref{ass:iid}--\ref{ass:construct}(a) and
\ref{ass:budget}--\ref{ass:sparsity},
\[
  \Delta_j:=\min_{\tau\in\cT_{\bar L}\setminus\cS_j}
  \Bigl\{R^{(j)}(\tau)-\min_{\tau'\in\cT_{\bar L}}R^{(j)}(\tau')\Bigr\}
  \ >\ 0 .
\]
\end{lemma}

The margin is a consequence of exact global optimisation over a finite class,
not a hypothesis: every $\tau\notin\cS_j$ has a leaf on which $\gamma_{0,j}$ is
non-constant with positive probability, hence strictly positive excess risk, and
a minimum of finitely many strictly positive numbers is strictly positive. The
two nuisances draw on different assumptions here. The propensity margin
$\Delta_e$ requires no positivity, since its excess risk is
$\Ltwo{\Pi_\tau\ez-\ez}^2$, strictly positive as soon as a leaf mixes two
positive-probability cells with distinct $\ez$. The outcome margin $\Delta_\mu$
requires Assumption~\ref{ass:causal}, because the outcome risk is
control-weighted and the relevant mass is $p_x\{1-\ez(x)\}\ge cp_x$.

\begin{proposition}[Greedy splitting has no margin]\label{prop:greedy}
There exist a finite product space $\cX$ and a propensity $\ez$ on it, with all
cells of positive probability and $\cS_e\ne\emptyset$ at $\bar L=4$, such that
the reduction in population risk achieved by every admissible first split is
exactly zero. Consequently the analogue of $\Delta_e$ for greedy recursive
partitioning is zero, while $\Delta_e>0$ by Lemma~\ref{lem:margin}.
\end{proposition}

\noindent Proposition~\ref{prop:greedy} is what converts Lemma~\ref{lem:margin}
from a statement about a finite class into a statement about the solver: the
margin is available to an exact optimiser and unavailable to a greedy one. The
witness, $\ez$ depending on $X$ only through $\mathbf1\{x_1=x_2\}$, is worked in
Appendix~\ref{app:counterex}. Assumption~\ref{ass:global} therefore cannot be
weakened to greedy recursive splitting, and this is the point at which the use
of globally optimal rather than greedily grown trees does mathematical rather
than expository work.

\subsection{Selection consistency}

\begin{lemma}[Selection lands in the sufficient class]\label{lem:selection}
Under Assumptions~\ref{ass:iid}--\ref{ass:global} and \ref{ass:sparsity}, for
each $j\in\{e,\mu\}$,
\[
  \PP\bigl(\that_j\in\cS_j\bigr)\longrightarrow1 .
\]
If in addition $Y$ is bounded, then
$\PP(\that_j\notin\cS_j)\le\card{\cT_{\bar L}}\,e^{-c_1n\Delta_j^2}$ for a
constant $c_1>0$ depending only on the bound and on $c$.
\end{lemma}

\noindent Only the exponential rate uses a tail condition. The convergence
$\PP(\that_j\in\cS_j)\to1$ holds under Assumption~\ref{ass:moments} alone,
because for each fixed $\tau$ the empirical risk converges by the law of large
numbers and uniformity over a fixed finite candidate set is automatic.

The penalty $\lambda_n$ in~\eqref{eq:select} is asymptotically inert for
inference: by Lemma~\ref{lem:invariance} every member of $\cS_j$ induces the
same estimand, so the penalty's role is to select the smallest sufficient tree,
that is, the most readily displayed one. Regularisation here serves
interpretability rather than statistical validity.

% ===========================================================================
\section{Full-sample inference with two optimal trees}\label{sec:main}

Both partitions and both sets of leaf values are computed from all $n$
observations. There is no sample splitting, no cross-fitting, no intersection of
candidate sets, and no averaging over fits.

\subsection{An exact decomposition}

Let $\varepsilon_i=Y_i-\mz(X_i)$, $\Delta\mu=\mhat-\mz$, $\Delta w=\what-\wz$,
and $\phi(O;\eta)=A\{Y-\mu(X)-\theta_0\}-(1-A)w(X)\{Y-\mu(X)\}$, so that
$\phi(\cdot;\eta_0)=\pi\psi_0$.

\begin{lemma}[Skeleton identity]\label{lem:skeleton}
On the event $\{\pihat>0\}$, and with no expansion or approximation,
\begin{equation}\label{eq:skeleton}
  \pihat\,(\thetahat-\theta_0)=\Pn\phi(\cdot;\eta_0)+U-T,
\end{equation}
where
\[
  U:=\Pn\bigl[h(O)\,\Delta\mu(X)\bigr],\quad
  h(O):=\wz(X)(1-A)-A,
  \qquad
  T:=\Pn\bigl[(1-A)\,\Delta w(X)\{Y-\mhat(X)\}\bigr],
\]
and $\E[h\mid X]=\wz(1-\ez)-\ez=0$ pointwise.
\end{lemma}

The two remainder terms have different characters and are controlled by
different arguments. Assumption~\ref{ass:sparsity} removes the first-order
approximation error; what survives in $U$ and $T$ is driven by
\emph{estimation} error of the leaf values, and is controlled by
Assumptions~\ref{ass:construct}--\ref{ass:budget}. Of the two, $U$ is the
binding constraint.

\begin{lemma}[Leaf-value term]\label{lem:U}
Under Assumptions~\ref{ass:iid}--\ref{ass:budget}, for every fixed
$\tau_\mu\in\cS_\mu$, $\abs U=O_p(\bar L/n)=o_p(n^{-1/2})$; the same bound holds
for the maximum over $\tau_\mu\in\cS_\mu$. The bound does not involve
$\tau_e$.
\end{lemma}

\begin{lemma}[Cross term]\label{lem:cross}
Under Assumptions~\ref{ass:iid}--\ref{ass:budget} and \ref{ass:sparsity},
\[
  \max_{(\tau_e,\tau_\mu)\in\cS_e\times\cS_\mu}
  \bigl\lvert T(\tau_e,\tau_\mu)\bigr\rvert
  =O_p\bigl(\sqrt{L_e}/n\bigr)=o_p(n^{-1/2}).
\]
\end{lemma}

\noindent Three features of Lemma~\ref{lem:cross} deserve emphasis, because they
determine the shape of its proof. First, the index set
$\cS_e\times\cS_\mu$ is tight: sufficiency of $\tau_\mu$ makes the substitution
$Y-\mhat\to\varepsilon$ exact, and sufficiency of $\tau_e$ is what permits
$\ez(X_i)$ to be collapsed to a single leaf value before the binomial variance
bound. Were $\tau_e\notin\cS_e$, then $\Delta w$ would converge to a non-zero
limit, centring within the leaves of $\tau_\mu$ would leave
$\sum_ia_i^2=\Theta_p(n)$, and $\sqrt nT$ would be $O_p(1)$ rather than
$o_p(1)$. Second, no minimum-leaf-mass condition is needed for this bound: each
leaf contributes at most a constant to the relevant sum irrespective of its
size, because the ratio of control count to total count in a leaf is at most one
deterministically. Third, clipping preserves the bound, being a projection onto
an interval containing the leaf value of the truth.

The cross term is not of product form. It pairs the propensity error with
outcome \emph{noise}, so the argument must exploit conditional mean-zeroness and
cannot exploit smallness. Cauchy--Schwarz is not merely lossy here but fatal:
its second factor converges to the control-arm residual standard deviation
$\E[\{1-\ez\}\sigma_0^2]^{1/2}$, which is bounded away from zero, leaving a
bound of order $\sqrt{L_e/n}\ge n^{-1/2}$ for every $L_e\ge1$. No leaf budget
rescues that, since the failure occurs at the scale of the central limit
theorem.

\subsection{Uniform asymptotic linearity}

For $(\tau_e,\tau_\mu)\in\cT_{\bar L}^2$ let $\thetahat(\tau_e,\tau_\mu)$ denote
the estimator~\eqref{eq:estimator} formed from the leaf refits on those
partitions.

\begin{lemma}[Uniform asymptotic linearity on the sufficient class]
\label{lem:uniform}
Under Assumptions~\ref{ass:iid}--\ref{ass:budget} and \ref{ass:sparsity},
\[
  \max_{(\tau_e,\tau_\mu)\in\cS_e\times\cS_\mu}
  \Bigl\lvert\sqrt n\bigl\{\thetahat(\tau_e,\tau_\mu)-\theta_0\bigr\}
   -\sqrt n\,\Pn\psi_0\Bigr\rvert=o_p(1).
\]
\end{lemma}

Lemma~\ref{lem:uniform} is what licenses the passage from fixed partitions to
the selected ones, and the form of the statement matters. Three routes that
appear adequate are not. Invariance of the \emph{estimand}
(Lemma~\ref{lem:invariance}) is invariance of first-order bias, whereas a
central limit theorem requires invariance of the \emph{influence function}. A
finite mixture $\sum_\tau\mathbf1\{\that=\tau\}\sqrt n\{\thetahat(\tau)-\theta_0\}$
is not asymptotically normal merely because each component is, since the
selection indicators are correlated with the components; finiteness of the
candidate set is not by itself sufficient, and what is needed is asymptotic
equivalence to a \emph{fixed} statistic. Conditioning on $\{\that_\mu=\tau\}$ is
not a repair, because leaf membership is determined by $X$ while the merging of
cells is determined by $Y$, so $\mhat$ remains conditionally biased given the
selection event.

Lemma~\ref{lem:uniform} conditions on nothing. It bounds a maximum over a family
that does not depend on the data, and a maximum of finitely many $o_p(n^{-1/2})$
terms is $o_p(n^{-1/2})$. It requires Assumption~\ref{ass:construct}(c) but
neither Assumption~\ref{ass:global} nor Assumption~\ref{ass:margin-floor}.

It is worth recording alongside that neither tree uses treated outcomes:
$\that_e$ and $\ehat$ are $\sigma(X_{1:n},A_{1:n})$-measurable, and $\that_\mu$
and $\mhat$ are $\sigma(X_{1:n},A_{1:n},\{Y_i:A_i=0\})$-measurable. The only
double use of the outcome is of the control arm, and the double use of the
treatment indicator is harmless because the propensity error enters
\eqref{eq:skeleton} only multiplied by a residual whose within-leaf control sum
is exactly zero.

\subsection{The main result}

\begin{theorem}[Full-sample two-tree central limit theorem]\label{thm:main}
Under Assumptions~\ref{ass:iid}--\ref{ass:global} and \ref{ass:sparsity}, with
$\that_e,\that_\mu$ selected by~\eqref{eq:select} on all $n$ observations and
leaf values refitted on all $n$ observations,
\[
  \sqrt n\bigl(\thetahat-\theta_0\bigr)\rightsquigarrow N(0,V),
\]
with $V$ as in~\eqref{eq:vdecomp}.
\end{theorem}

\begin{corollary}[Variance estimation and confidence intervals]
\label{cor:variance}
Under the conditions of Theorem~\ref{thm:main}, let
\begin{equation}\label{eq:vhat}
  \Vhat=\frac1{n\pihat^2}\sum_{i=1}^n
  \Bigl[A_i\bigl\{Y_i-\mhat(X_i)-\thetahat\bigr\}
   -(1-A_i)\what(X_i)\bigl\{Y_i-\mhat(X_i)\bigr\}\Bigr]^2 .
\end{equation}
Then $\Vhat\to_pV$, and $\thetahat\pm z_{1-\alpha/2}\sqrt{\Vhat/n}$ has
asymptotic coverage $1-\alpha$ for $\theta_0$, pointwise in $P$.
\end{corollary}

\noindent The route to Corollary~\ref{cor:variance} is a union bound over the
finite sufficient class, and it requires three inputs:
Lemma~\ref{lem:invariance}, so that every $\tau\in\cS_j$ has the same limit and
$\plim\Vhat$ is not a selection-weighted mixture of distinct numbers;
Lemma~\ref{lem:selection}, to discard the complementary event; and
Lemma~\ref{lem:findim} below, for consistency at fixed $\tau$. Notably it
requires only $\Linf{\widehat\eta-\eta_0}=o_p(1)$, and neither the rate
condition of Section~\ref{sec:crossfit} nor Lemma~\ref{lem:cross}.

\begin{lemma}[Conditional finite-dimensionality]\label{lem:findim}
Fix $\tau\in\cT_{\bar L}$ all of whose leaves have positive probability under
$P$, and, for $j=\mu$, positive control mass. Then the leaf refit for nuisance
$j$ is determined by at most $\card\tau$ real numbers, and if $\tau\in\cS_j$,
\[
  \max_{\ell\in\tau}
  \bigl\lvert\widehat\gamma_{j}\restriction_\ell
   -\gamma_{0,j}\restriction_\ell\bigr\rvert=O_p(n^{-1/2}).
\]
\end{lemma}

Two finite-sample properties of~\eqref{eq:vhat} are worth stating. Its summands
have exactly zero sample mean by construction, so it is simultaneously the
uncentred second moment and the centred sample variance, and no centring
ambiguity arises. Because structure and leaf values are both fitted on all $n$
observations, the summands are in-sample residuals, and~\eqref{eq:vhat} is
downward biased by a relative factor of order $(2\bar L+1)/n$; a
degrees-of-freedom divisor $n-(2\card{\that_e}+\card{\that_\mu})$ is a
reasonable finite-sample convention. All statements are on $\{\pihat>0\}$, whose
probability tends to one.

\subsection{Uniformity in the underlying law}\label{sec:uniformity}

Lemma~\ref{lem:selection} gives
$\PP(\that_j\notin\cS_j)\le\card{\cT_{\bar L}}e^{-c_1n\Delta_j^2}$, which is
$o(1)$ for each fixed $P$ but not uniformly over
$\{P:\delta_e=\delta_\mu=0\}$, since $\Delta_j$ may be arbitrarily small in that
set. Coverage in Corollary~\ref{cor:variance} is therefore asserted pointwise.
Uniformity is recovered on a class that fixes the margin together with the
moment, positivity, and non-degeneracy constants.

\begin{definition}[Uniformity class]\label{def:uniformity}
For $\Delta_{\min},c,v_0>0$ and $C_4<\infty$, let
\[
  \cP(\Delta_{\min},c,C_4,v_0)
  =\Bigl\{P:\ \delta_e=\delta_\mu=0,\ \min_j\Delta_j\ge\Delta_{\min},\
  1-\ez\ge c,\ \E[Y^4\mid A,X]\le C_4,\ V\ge v_0\Bigr\}.
\]
\end{definition}

\begin{theorem}[Uniform coverage]\label{thm:uniform-cover}
Under Assumptions~\ref{ass:construct}--\ref{ass:global}, the interval
$\thetahat\pm z_{1-\alpha/2}\sqrt{\Vhat/n}$ has coverage converging to
$1-\alpha$ uniformly over $\cP(\Delta_{\min},c,C_4,v_0)$.
\end{theorem}

\noindent All four constants are needed. Without the moment and non-degeneracy
constraints $V$ may be unbounded or $V>0$ may fail, so uniform coverage cannot
hold; and uniform-in-$P$ selection consistency requires a fourth moment, which
Assumption~\ref{ass:moments} does not supply. No tail condition is needed for
the pointwise results; the uniform version requires a uniform fourth moment, or
a bounded outcome.

\subsection{Regularity and efficiency}\label{sec:regularity}

\begin{theorem}[Regularity and local asymptotic efficiency]\label{thm:regular}
Under the conditions of Theorem~\ref{thm:main}, $\thetahat$ is regular at $P_0$
relative to $\cP$: for every path $\{P_{n,h}\}$ in $\cP$ that is differentiable
in quadratic mean at $P_0$ with score $h$,
\[
  \sqrt n\bigl\{\thetahat-\theta(P_{n,h})\bigr\}\rightsquigarrow N(0,V)
  \qquad\text{under }P_{n,h}^n,
\]
with the limit free of $h$. Consequently, by the convolution theorem,
$\thetahat$ is locally asymptotically efficient in $\cP$.
\end{theorem}

The efficiency claim is asserted in $\cP$ and must not be asserted in the
sparsity-restricted submodel, where the bound is strictly smaller.

\begin{proposition}[The bound in the sparse submodel]\label{prop:sparse-bound}
For $\tau\in\cS_e$, let $\cP_\tau\subset\cP$ impose that $\ez$ be constant on
the leaves of $\tau$. The efficiency bound for $\theta_0$ in $\cP_\tau$ is
\[
  V_\tau=V-\pi^{-2}\,\E\Bigl[\ez(1-\ez)\bigl\{\Delta m
   -\Pi^{v}_\tau\Delta m\bigr\}^2\Bigr]\ \le\ V,
  \qquad \Delta m=m_1-\mz-\theta_0,\quad v_x=p_x\ez(x)\{1-\ez(x)\},
\]
with strict inequality whenever the treated--control contrast varies within a
leaf of $\tau$. Moreover $V^{\mathrm{known}\text{-}e}\le V_\tau\le V$, where the
first is the bound when $\ez$ is known.
\end{proposition}

\noindent For \emph{known} $\tau$ the bound $V_\tau$ is attainable by a regular
estimator; this is the classical asymmetry that for the ATT, knowledge of the
propensity lowers the bound \citep{hahn1998}. The obstruction is that $\tau$ is
unknown. The relevant quantity across the union of sparse submodels is
$V_{\mathrm{sparse}}=\max_{\tau\in\cS_e}V_\tau$, governed by the finest
sufficient partition within the budget, and hence by $\bar L$ as well as by $P$;
it depends discontinuously on $P$, because $\cS_e$ can empty out as $P$ leaves
the sparse class. Accordingly no estimator attains $V_{\tau(P)}$ at every
$P\in\bigcup_\tau\cP_\tau$ while remaining regular. The estimator studied here
is regular, is uniformly valid over the class of
Definition~\ref{def:uniformity}, and attains $V$.

% ===========================================================================
\section{Cross-fitted inference with optimal-tree nuisances}\label{sec:crossfit}

The estimator of this section drops Assumption~\ref{ass:sparsity} and imposes no
leaf budget, at the cost of $K$ trees per nuisance rather than one. It is
included for two reasons: it is what remains available when structural sparsity
is untenable, and it is the anchor for Section~\ref{sec:honest}.

Partition $\{1,\dots,n\}$ into $K$ folds of size $n_k$ with $K$ fixed. For each
$k$, fit $\widehat\eta^{(-k)}$ by~\eqref{eq:select} and~\eqref{eq:refit} on the
observations outside fold $k$, and let $\thetahat_{\mathrm{cf}}$ solve
$\sum_k\sum_{i\in\text{fold }k}\psi(O_i;\theta,\widehat\eta^{(-k)})=0$.

\begin{assumption}[Rate condition]\label{ass:rate}
$\Ltwo{\ehat-\ez}\cdot\Lw{\mhat-\mz}=o_p(n^{-1/2})$, where the norms are
evaluated at the fold-wise fits.
\end{assumption}

\begin{theorem}[Cross-fitted central limit theorem]\label{thm:crossfit}
Under Assumptions~\ref{ass:iid}--\ref{ass:construct}, \ref{ass:global},
and~\ref{ass:rate}, with $K$-fold cross-fitting and $K$ fixed,
\[
  \sqrt n\bigl(\thetahat_{\mathrm{cf}}-\theta_0\bigr)\rightsquigarrow N(0,V).
\]
\end{theorem}

\noindent Theorem~\ref{thm:crossfit} is an application of the general theory of
debiased machine learning \citep{chernozhukov2018}: Neyman orthogonality of
\eqref{eq:score} is immediate from Proposition~\ref{prop:bilinear}, and the
product structure of~\eqref{eq:convert} is what Assumption~\ref{ass:rate}
addresses. The pairing in Assumption~\ref{ass:rate} is unweighted for the
propensity and control-weighted for the outcome, exactly as
Corollary~\ref{cor:convert} delivers.

Assumption~\ref{ass:rate} is verifiable rather than merely assumable in the
present setting.

\begin{proposition}[The rate condition holds at the saturated budget]
\label{prop:rate-verify}
Under Assumptions~\ref{ass:iid}--\ref{ass:global} with $\bar L=M$,
Assumption~\ref{ass:sparsity} holds automatically and
Assumption~\ref{ass:rate} is satisfied with
$\Ltwo{\ehat-\ez}\cdot\Lw{\mhat-\mz}=O_p(n^{-1})$.
\end{proposition}

\noindent The argument is short: at $\bar L=M$ the singleton partition is
tree-realisable using $M-1$ splits, so $\cS_e$ and $\cS_\mu$ are non-empty and
Assumption~\ref{ass:sparsity} carries no content; Lemma~\ref{lem:margin} then
applies, and Lemma~\ref{lem:selection} places $\that_j$ in $\cS_j$ with
probability tending to one; on that event $\gamma_{0,j}$ is leaf-constant, so
each fitted nuisance is a leaf mean of a leaf-constant truth and converges at
$O_p(\sqrt{M/n})=O_p(n^{-1/2})$ by Lemma~\ref{lem:findim}; the product is
$O_p(n^{-1})$.

The appeal to Lemma~\ref{lem:margin} in Proposition~\ref{prop:rate-verify} is
not dispensable. The alternative route, which uses only finiteness of the
candidate class and a union bound without a margin, yields excess risk
$O_p(n^{-1/2})$ per nuisance, hence $O_p(n^{-1/4})$ in norm, hence a product of
exactly $O_p(n^{-1/2})$ --- which fails Assumption~\ref{ass:rate}, since the
requirement is $o_p(n^{-1/2})$. Finiteness of the candidate set is not by itself
enough.

% ===========================================================================
\section{Inference when sparsity fails}\label{sec:honest}

Without Assumption~\ref{ass:sparsity}, Lemma~\ref{lem:invariance} fails and
Lemma~\ref{lem:biasbound} becomes the object to control. Validity survives;
efficiency does not; and the price is estimable. Two constructions are given.
The first requires no assumption beyond a valid anchor and is the recommended
default. The second is sharper but requires
Assumption~\ref{ass:pseudo} and a conservative variance estimator.

\subsection{The anchor interval}

\begin{theorem}[Anchor interval]\label{thm:anchor}
Let $\thetatil$ satisfy
$\sqrt n(\thetatil-\theta_0)\rightsquigarrow N(0,\widetilde V)$ with
$\widetilde\sigma^2\to_p\widetilde V$, and let
$\widehat\delta=\thetahat-\thetatil$. Then
\[
  C_n=\Bigl[\thetahat\pm\bigl(\abs{\widehat\delta}
   +z_{1-\alpha/2}\,\widetilde\sigma/\sqrt n\bigr)\Bigr]
\]
satisfies $\liminf_n\PP(\theta_0\in C_n)\ge1-\alpha$, with no condition on
$\delta_e$, $\delta_\mu$, $\Dw$, or $\Dm$, and no smoothness condition on the
nuisances.
\end{theorem}

\noindent Theorem~\ref{thm:crossfit} supplies the anchor. The interval $C_n$
contains the anchor's own interval recentred at $\thetahat$ and is therefore
never narrower than the cross-fitted interval; what it buys is that the point
estimate remains the one the displayed trees produce. The construction is
deliberately crude, requiring nothing beyond the triangle inequality; a
bias-aware critical value \citep{armstrong2018}, which exploits the fact that a
bias of known magnitude $B$ costs less than $B$ in half-width, is preferable in
practice.

The difference $\widehat\delta$ must be used to widen the interval and never to
correct the point estimate, since $\thetahat-\widehat\delta=\thetatil$
identically: subtracting it returns the cross-fitted estimator exactly and
forfeits the displayed-tree estimate for nothing.

\begin{corollary}[Width]\label{cor:width}
Under the conditions of Theorem~\ref{thm:anchor},
$\abs{\widehat\delta}=O_p(\delta_e\delta_\mu+n^{-1/2})$ and
\[
  \mathrm{width}(C_n)\asymp_p\max\bigl\{n^{-1/2},\ \delta_e\delta_\mu\bigr\}.
\]
The interval contracts at the parametric rate if and only if
$\delta_e\delta_\mu=O(n^{-1/2})$, and its width otherwise plateaus at the bias
floor.
\end{corollary}

\noindent Corollary~\ref{cor:width} quantifies the trade-off. The correct
statement is not that a bounded-complexity tree cannot deliver $\sqrt n$
inference for $\theta_0$ --- it can, by Theorem~\ref{thm:main} --- but that the
cost of insisting on a displayable tree is an interval whose width contracts at
the parametric rate only while the tree is rich enough to track the nuisances,
and plateaus thereafter at a height equal to the product of the approximation
errors. With $\bar L$ fixed that height does not shrink.

\subsection{The fidelity diagnostic}

\begin{proposition}[Fidelity diagnostic]\label{prop:spectest}
Under Assumptions~\ref{ass:iid}--\ref{ass:global}, \ref{ass:sparsity},
and~\ref{ass:rate}, $\sqrt n\,\widehat\delta=o_p(1)$. Consequently, for any
$\widehat{\mathrm{SE}}\asymp n^{-1/2}$, the test that rejects
Assumption~\ref{ass:sparsity} when
$\abs{\widehat\delta}>z_{1-\alpha/2}\widehat{\mathrm{SE}}$ has size tending to
zero and is consistent against fixed alternatives.
\end{proposition}

\noindent Under Assumption~\ref{ass:sparsity} both estimators are
asymptotically linear with the \emph{same} influence function $\psi_0$, by
Theorem~\ref{thm:main} and Theorem~\ref{thm:crossfit}, so the leading terms of
$\widehat\delta$ cancel identically. The diagnostic is therefore a conservative
test rather than an asymptotically exact one, and it admits no non-degenerate
limiting variance. It cannot serve as a widening term on its own: it is
$O_p(n^{-1/2})$ even when the bias is exactly zero, so it neither centres on nor
dominates the bias.

\subsection{The realised-error refinement}\label{sec:realised}

The anchor interval pays the anchor's full width. A sharper interval is
available under Assumption~\ref{ass:pseudo}, which supplies the limits that
misspecification otherwise destroys: without a population margin the minimiser
in~\eqref{eq:select} may oscillate among near-ties, in which case
$\widehat\eta$ has no limit and $V_*$ is not defined.

\begin{proposition}[Realised errors are estimable]\label{prop:estimable}
Suppose Assumptions~\ref{ass:iid}--\ref{ass:budget} and~\ref{ass:pseudo} hold,
$Y$ is bounded or sub-Gaussian conditionally on $(A,X)$, every cell is occupied
in the relevant arm in the sense that $np_{\min}\gtrsim\log M$ for $e$ and
$ncp_{\min}\gtrsim\log M$ for $\mu$, and the saturated cell-wise fits
$\ehat_{\mathrm{sat}},\mhat_{\mathrm{sat}}$ are computed on a sample independent
of that used for $\widehat\eta$. Then the degrees-of-freedom-corrected
statistic
\[
  \widehat D_\mu^2=\sum_x\widehat p_x\{1-\ehat(x)\}
   \bigl\{\mhat(x)-\mhat_{\mathrm{sat}}(x)\bigr\}^2
  -\sum_x\widehat p_x\{1-\ehat(x)\}\frac{\widehat\sigma^2(x)}{n_{0,x}},
\]
and its analogue $\widehat D_w^2$, satisfy
$\widehat D_j^2-D_j^2=O_p\bigl(M/n+D_jn^{-1/2}\bigr)$ up to logarithmic factors.
Setting $U_j^2=(\widehat D_j^2)_++\kappa_na_{n,j}$ with
$a_{n,j}=(M\log n)/n+\widehat D_j\sqrt{\log n/n}$ and $\kappa_n\to\infty$ at a
polylogarithmic rate yields
$\PP\bigl(U_w\ge\Dw\text{ and }U_\mu\ge\Dm\bigr)\to1$.
\end{proposition}

\begin{theorem}[Interval under a pseudo-true margin]\label{thm:honest}
Suppose Assumptions~\ref{ass:iid}--\ref{ass:budget} and~\ref{ass:pseudo} hold,
let $U_w,U_\mu$ satisfy
$\PP(U_w\ge\Dw,\,U_\mu\ge\Dm)\to1$, and let $\Vhat^{+}$ satisfy
$\plim\Vhat^{+}\ge V_*$. Then
\[
  \thetahat\ \pm\ \Bigl(z_{1-\alpha/4}\sqrt{\Vhat^{+}/n}
   +\frac{U_wU_\mu}{\pi}\Bigr)
\]
has asymptotic coverage at least $1-\alpha$ for $\theta_0$.
\end{theorem}

\noindent Two features of Theorem~\ref{thm:honest} are not interchangeable with
their analogues in Theorem~\ref{thm:main}. The comparison variance is $V_*$, the
asymptotic variance at the pseudo-true nuisances, and not $V$: under
misspecification the terms $U$ and $T$ of~\eqref{eq:skeleton} contribute
variance absent from $V$, so $V_*$ may exceed $V$ and a bound on $V$ does not
bound $V_*$. The plug-in~\eqref{eq:vhat} does not qualify as $\Vhat^{+}$: with
$a=\mu_*-\mz$ and $b=w_*-\wz$,
\[
  \plim\Vhat=V+2\pi^{-2}\bigl\{\E[(1-\ez)\wz b\sigma_0^2]
   -\E[\ez\,a\,\Delta m]\bigr\}+O\bigl(\Ltwo a^2+\Ltwo b^2\bigr),
\]
whose first-order term is not sign-definite; taking $b\equiv0$ and
$a=\epsilon(1-\ez)\Delta m$ makes it negative against a second-order
$O(\epsilon^2)$, so $\plim\Vhat<V$ is possible. The error budget in
Theorem~\ref{thm:honest} allocates $\alpha/2$ to the normal approximation and
$\alpha/4$ to each one-sided bound, summing to $\alpha$.

Proposition~\ref{prop:estimable} requires a hold-out for the saturated
reference, which is in tension with the no-splitting design of
Sections~\ref{sec:main} and~\ref{sec:honest}; the requirement is stated in the
proposition because it cannot be removed without a separate argument for the
dependence between $\mhat$ and $\mhat_{\mathrm{sat}}$, which share the control
outcomes. Its rates are stated in $M$, and are informative only when $M$ is
regarded as large relative to $\bar L$ rather than as the fixed constant of
Assumption~\ref{ass:finite}; read literally under
Assumption~\ref{ass:finite}, $M/n=O(n^{-1})$ and the conditions are vacuous.
Proposition~\ref{prop:estimable} is accordingly a finite-sample guide to when
the refinement is worth its cost, and Theorem~\ref{thm:anchor} is the
asymptotic statement.

% ===========================================================================
\section{Discussion}\label{sec:discussion}

\paragraph{What the results establish.}
Three claims are available, and they answer different questions. Inside the
sparse class, Theorem~\ref{thm:main} and Corollary~\ref{cor:variance} give
efficient $\sqrt n$ inference for $\theta_0$ from the displayed trees, with the
ordinary plug-in variance and no sample splitting. Outside it with the bias
bounded, Theorem~\ref{thm:anchor} and Corollary~\ref{cor:width} give valid but
inefficient inference at width $\max\{n^{-1/2},\delta_e\delta_\mu\}$, the point
estimate still being the one the trees produce. Outside it with the bias
unbounded, the trees remain an audit object reported alongside the cross-fitted
estimate, with $\widehat\delta$ as a diagnostic and no inferential claim
attached to the trees.

\paragraph{Interpretability.}
The displayed model is the estimator. It is neither a post hoc surrogate for a
black-box nuisance nor a collection of $K$ fold-specific trees: one propensity
tree and one outcome tree, fitted on every observation, are simultaneously the
audit object and the inferential object. The regularisation that makes them
small is statistically inert by Lemma~\ref{lem:invariance}, so
interpretability is obtained at no cost in validity.

\paragraph{A shared partition degenerates.}
On the common refinement $\that_e\wedge\that_\mu$ both $\what$ and $\mhat$ are
leaf-constant, so the identity $\sum_{i\in\ell,A_i=0}(Y_i-\mhat_\ell)=0$ makes
the control-arm term of~\eqref{eq:estimator} identically zero and the estimator
reduces to the treated-weighted stratified difference in means
$\sum_\ell(n_{1,\ell}/n_1)(\bar Y_{1,\ell}-\bar Y_{0,\ell})$, an exact
finite-sample identity unaffected by clipping. That estimator is consistent and
asymptotically normal with $T=0$ identically, so its analysis is easier. The
reasons not to prefer it are not reasons of validity: a single displayed tree
would carry no fitted propensity values, defeating the purpose of displaying it;
and $\card{\that_e\wedge\that_\mu}$ can be as large as $L_eL_\mu$, so the
refinement may violate Assumption~\ref{ass:construct}(d) and
Assumption~\ref{ass:budget} even when each tree satisfies them. The propensity's
fitted values are annihilated in the common refinement, but its structure still
selects the stratification and so determines both the point estimate and the
variance.

\paragraph{Boundaries.}
Three extensions are outside the scope of this document. Growing complexity,
either $M_n\to\infty$ or $\bar L_n\to\infty$, makes $\cT_{\bar L}$
non-fixed, so that uniformity is no longer free and the entropy
$\bar L_n\log(dJ_n)$ becomes active together with a tail condition; the order in
which the present arguments fail as $\bar L_n$ grows is
$U=O_p(\bar L_n/n)$ first, at $\bar L_n\asymp\sqrt n$, then
Lemma~\ref{lem:selection}, then Lemma~\ref{lem:cross}, which merely reinstates a
logarithmic factor, and Lemma~\ref{lem:uniform} last, which would require
genuine localisation. A sufficient set of conditions for that regime is
$\bar L_n=o(\sqrt n)$, $\bar L_n\log(d_nJ_n)=o(n\Delta_{\min,n}^2)$, and a
bounded or sub-Gaussian outcome. Continuous covariates require approximation
theory and an oracle inequality for optimal trees, and Rashomon sets require
structure-intersection results; both are developed in the companion paper and
neither is used above. A sharp conservative variance estimator $\Vhat^{+}$ for
Theorem~\ref{thm:honest}, as opposed to the crude bound available from
Assumptions~\ref{ass:causal} and~\ref{ass:moments}, is not constructed here.

\paragraph{Assumed decay as an alternative to measurement.}
The alternative to estimating the approximation errors is to assume their decay,
$\delta_j(L)\le C_jL^{-q_j}$ with $q_e+q_\mu>1$, the discrete counterpart of
assuming a Hölder or Besov ball. Weak-$\ell_r$ decay of the tree coefficients
gives $q=1/r-1/2$. No decay rate follows from finiteness alone: parity on $d$
binary coordinates has $\delta(L)$ bounded away from zero for every
axis-aligned tree with $L<2^{d-1}$ leaves and zero thereafter, so decay may be a
step function of $L$. Decay of low-order interactions is not an appropriate
sufficient condition here, since such functions are well approximated by
additive models, which are the least favourable case for trees.

% ===========================================================================
\appendix

\section{Auxiliary results}\label{app:aux}

\emph{To be completed: norm-equivalence bounds, the non-expansiveness of
one-sided clipping, and the leaf-mean moment bounds used repeatedly in
Appendices~\ref{app:structure} and~\ref{app:main}.}

\section{Proofs for Section~\ref{sec:eif}}\label{app:eif}

\emph{To be completed: Proposition~\ref{prop:bilinear},
Lemma~\ref{lem:biasbound}, Corollary~\ref{cor:convert},
Lemma~\ref{lem:tangent}.}

\section{Proofs for Section~\ref{sec:structure}}\label{app:structure}

\emph{To be completed: Proposition~\ref{prop:pinwheel},
Lemma~\ref{lem:invariance}, Lemma~\ref{lem:margin},
Proposition~\ref{prop:greedy}, Lemma~\ref{lem:selection}.}

\section{Proofs for Section~\ref{sec:main}}\label{app:main}

\emph{To be completed: Lemma~\ref{lem:skeleton}, Lemma~\ref{lem:U},
Lemma~\ref{lem:cross}, Lemma~\ref{lem:uniform}, Lemma~\ref{lem:findim},
Theorem~\ref{thm:main}, Corollary~\ref{cor:variance},
Theorem~\ref{thm:uniform-cover}, Theorem~\ref{thm:regular},
Proposition~\ref{prop:sparse-bound}.}

\section{Proofs for Sections~\ref{sec:crossfit} and~\ref{sec:honest}}
\label{app:honest}

\emph{To be completed: Theorem~\ref{thm:crossfit},
Proposition~\ref{prop:rate-verify}, Theorem~\ref{thm:anchor},
Corollary~\ref{cor:width}, Proposition~\ref{prop:spectest},
Proposition~\ref{prop:estimable}, Theorem~\ref{thm:honest}.}

\section{Counterexamples}\label{app:counterex}

\emph{To be completed: the five-box pinwheel of
Proposition~\ref{prop:pinwheel}; the zero-margin greedy witness of
Proposition~\ref{prop:greedy}; and the borderline rate computation following
Proposition~\ref{prop:rate-verify}.}

\section{Crosswalk to the submitted paper}\label{app:crosswalk}

\begin{center}
\begin{longtable}{@{}llp{4.6cm}c@{}}
\toprule
Result & Statement & Paper location & Pages \\ \midrule
\endfirsthead
\toprule
Result & Statement & Paper location & Pages \\ \midrule
\endhead
\multicolumn{4}{@{}l}{\emph{Main text}}\\[2pt]
--- & Setup, target, score \eqref{eq:score}
  & §2 Setting and target & 1.5 \\
A\ref{ass:iid}--A\ref{ass:global}
  & Sampling model and estimator & §2.1--2.2, displayed block & 1 \\
Def.~\ref{def:sufficient} & Sufficient class & §3.1 & 0.3 \\
A\ref{ass:sparsity} & Structural sparsity & §3.1, with the asymmetry discussion & 0.7 \\
Lem.~\ref{lem:invariance} & Invariance of the estimand & §3.2, proof inline & 0.5 \\
Lem.~\ref{lem:margin} & Derived margin & §3.3, proof inline & 0.4 \\
Prop.~\ref{prop:greedy} & Greedy has no margin & §3.3, witness in Appendix & 0.3 \\
Lem.~\ref{lem:selection} & Selection consistency & §3.4 & 0.3 \\
Lem.~\ref{lem:skeleton} & Skeleton identity & §4.1, displayed & 0.4 \\
Lem.~\ref{lem:cross} & Cross term & §4.2, statement only & 0.4 \\
Lem.~\ref{lem:uniform} & Uniform asymptotic linearity & §4.2, statement only & 0.4 \\
Thm.~\ref{thm:main} & Main central limit theorem & §4.3 & 0.5 \\
Cor.~\ref{cor:variance} & Variance and intervals & §4.3 & 0.4 \\
Thm.~\ref{thm:uniform-cover} & Uniform coverage & §4.4 & 0.3 \\
Thm.~\ref{thm:regular} & Regularity and efficiency & §4.5 & 0.5 \\
Prop.~\ref{prop:sparse-bound} & Bound in the sparse submodel & §4.5, condensed & 0.4 \\
Thm.~\ref{thm:crossfit} & Cross-fitted limit theorem & §5.1 & 0.4 \\
Prop.~\ref{prop:rate-verify} & Rate condition verified & §5.2 & 0.3 \\
Thm.~\ref{thm:anchor} & Anchor interval & §6.1, proof inline & 0.4 \\
Cor.~\ref{cor:width} & Width and the plateau & §6.1 & 0.3 \\
Prop.~\ref{prop:spectest} & Fidelity diagnostic & §6.2 & 0.3 \\
Prop.~\ref{prop:estimable}, Thm.~\ref{thm:honest}
  & Realised-error refinement & §6.3, condensed to one paragraph
  and a display & 0.5 \\
--- & Simulation design and results & §7 & 3 \\
--- & Facility-quality application & §8 & 2.5 \\
--- & Discussion & §9 & 1.5 \\[4pt]
\multicolumn{4}{@{}l}{\emph{Supplementary material}}\\[2pt]
App.~\ref{app:aux}--\ref{app:honest} & All proofs & Supplement §S1--S5 & --- \\
App.~\ref{app:counterex} & Counterexamples & Supplement §S6 & --- \\
--- & Continuous covariates, Rashomon sets & Companion paper & --- \\
\bottomrule
\end{longtable}
\end{center}

\begin{thebibliography}{9}
\bibitem[Armstrong \& Kolesár(2018)]{armstrong2018} Armstrong, T. B. and
  Kolesár, M. (2018). Optimal inference in a class of regression models.
  \emph{Econometrica} 86, 655--683.
\bibitem[Chen et al.(2022)]{chen2022} Chen, Q., Syrgkanis, V. and Austern, M.
  (2022). Debiased machine learning without sample splitting.
  arXiv:2206.01825.
\bibitem[Chernozhukov et al.(2018)]{chernozhukov2018} Chernozhukov, V.,
  Chetverikov, D., Demirer, M., Duflo, E., Hansen, C., Newey, W. and Robins, J.
  (2018). Double/debiased machine learning for treatment and structural
  parameters. \emph{The Econometrics Journal} 21, C1--C68.
\bibitem[Hahn(1998)]{hahn1998} Hahn, J. (1998). On the role of the propensity
  score in efficient semiparametric estimation of average treatment effects.
  \emph{Econometrica} 66, 315--331.
\bibitem[Kennedy(2024)]{kennedy2024} Kennedy, E. H. (2024). Semiparametric
  doubly robust targeted double machine learning: a review.
\bibitem[Leeb \& Pötscher(2005)]{leeb2005} Leeb, H. and Pötscher, B. M. (2005).
  Model selection and inference: facts and fiction. \emph{Econometric Theory}
  21, 21--59.
\bibitem[Robins \& Rotnitzky(1995)]{robins1995} Robins, J. M. and Rotnitzky, A.
  (1995). Semiparametric efficiency in multivariate regression models with
  missing data. \emph{Journal of the American Statistical Association} 90,
  122--129.
\bibitem[van der Laan \& Robins(2003)]{vanderlaan2003} van der Laan, M. J. and
  Robins, J. M. (2003). \emph{Unified Methods for Censored Longitudinal Data
  and Causality}. Springer.
\bibitem[van der Vaart(1998)]{vandervaart1998} van der Vaart, A. W. (1998).
  \emph{Asymptotic Statistics}. Cambridge University Press.
\end{thebibliography}

\end{document}
